\documentclass[11pt,letterpaper]{article}
\pdfinfoomitdate=1
\pdftrailerid{}
\pdfsuppressptexinfo=15
\usepackage[T1]{fontenc}
\usepackage[utf8]{inputenc}
\usepackage{lmodern,amsmath,amssymb,amsthm}
\usepackage[margin=1in]{geometry}
\usepackage{booktabs}
\usepackage[mathlines]{lineno}
\usepackage[hidelinks]{hyperref}
\hypersetup{pdftitle={Reflected Blocks, Implicit Parameters, and Full Geometric Recovery},
 pdfauthor={Maurizio Falconi}}
\newtheorem{theorem}{Theorem}[section]
\newtheorem{lemma}[theorem]{Lemma}
\newtheorem{corollary}[theorem]{Corollary}
\newcommand{\Rstar}{R^*}
\newcommand{\thR}{\theta_R}
\newcommand{\dd}{\,\mathrm d}
\newcommand{\al}{\widehat\alpha}
\newcommand{\epsb}{\epsilon_b}
\newcommand{\Ds}{\Delta_*}
\newcommand{\norm}[1]{\left\lVert#1\right\rVert}
\title{Reflected Blocks, Implicit Parameters,\\and Full Geometric Recovery}
\author{Maurizio Falconi}
\date{24 October 2026}
\begin{document}
\maketitle
\linenumbers
\begin{abstract}
This supplement gives a complete proof of the reflected-block upper
construction and of the implicit parameters defining
$U_4=C_4(1/20000)$. It includes the all-pairs angular criterion, recovery
by actual integer labels, a quantitative radius estimate, uniqueness of
the successive parameters, and the rational gates used to isolate them.
The universal statements are analytic theorems; the accompanying bounded
rational computations establish specified constant inequalities. The
construction is an upper bound and does not assert an optimal partition
or a finite geometric optimum. These are existing results reorganized
from the source proofs listed in the dependency document.
\end{abstract}
\tableofcontents

\section{The Geometric Problem and the Target Bound}
Circles of radii $1,\ldots,n$ are externally tangent to a central circle
of radius $R>0$. Let $\Rstar(n)$ be the infimum of feasible central radii.
For a prescribed cyclic order $\sigma$, let $R_{\rm full}(\sigma)$ use
all pairs and let $R_{\rm chain}(\sigma)$ use only adjacent gaps. The
cosine law gives the required angle
\begin{equation}\label{eq:angle}
 \thR(a,b)=2\arcsin\sqrt{\frac{ab}{(R+a)(R+b)}}.
\end{equation}
Both directed arcs between each pair must be at least this angle. This is
equivalent to Cartesian non-overlap, since their sum is $2\pi$ and the
shorter angle lies in $[0,\pi]$. Throughout, inverse trigonometric
functions have their principal real branches. The kernel increases in
each outer radius and decreases strictly in $R$. In general
$R_{\rm chain}\le R_{\rm full}$; no equality is presumed here.

Define, on $\mathcal D=[0,1]\times[1,2]^2$,
\begin{equation}\label{eq:g}
 g(t,X,Y)=\max\{\sqrt t(\sqrt X+\sqrt Y),\sqrt{XY}\},\qquad
 C(\mu)=\frac1{4\pi}\int g\dd\mu.
\end{equation}
Sections~\ref{sec:x}--\ref{sec:widths} uniquely define $x_*,\al,
\epsb,\Ds$. Set $A=1+\al$, $\lambda=Ax_*$,
$v=\lambda+\epsb$, and $w=v+\Ds$. The measure $\mu_4(\eta)$ consists
of four separate symmetric reflections on the consecutive intervals of
lengths $(\lambda,\epsb,\Ds,\eta)$, followed by the shifted diagonal;
its precise general definition is \eqref{eq:mu}. Put
$C_3(d)=C(\mu_{\al;(\lambda,\epsb,d)})$ and
$C_4(\eta)=C(\mu_4(\eta))$.
The final result proved here is
\begin{equation}\label{eq:target}
 \limsup_{n\to\infty}\frac{\Rstar(n)}{n^2}
 \le U_4:=C_4(1/20000)
 <C_3(\Ds)-\frac1{4608000000000000}.
\end{equation}
If $C_*$ denotes the global limit established in the main article,
the left side is $C_*$. Its existence is not needed for this supplement's
upper-bound proof. Terminating decimal bounds below are exact rationals,
never definitions of the implicit parameters.

\section{An Exact All-Pairs Criterion}
\begin{lemma}[High-Shell Inequality]\label{lem:shell}
If $0<u\le v\le2u$, then $2\thR(u,u)>\thR(v,v)$.
Consequently every path $H_0,\ldots,H_j$ in $[u,v]$, $j\ge1$, satisfies
$\sum_{i=1}^j\thR(H_{i-1},H_i)\ge\thR(H_0,H_j)$.
\end{lemma}
\begin{proof}
Write $x=u/(R+u)$ and $z=v/(R+v)\le2x/(1+x)$. If $x\ge1/\sqrt2$,
then $2\arcsin x\ge\pi/2>\arcsin z$. Otherwise
\[
 (1+x)^2(1-x^2)-1=x(2-2x^2-x^3)>0,
\]
because $x^2<1/2$, $x<3/4$, and $2-2x^2-x^3>5/8$.
Thus $z<2x\sqrt{1-x^2}=\sin(2\arcsin x)$. Both comparison angles
are in $(0,\pi/2)$, so the same strict inequality follows. Multiplying
by two proves the first assertion. For any $a,b,c\in[u,v]$,
\[
 \thR(a,b)+\thR(b,c)\ge2\thR(u,u)>\thR(v,v)\ge\thR(a,c).
\]
Induction contracts a path. The one-edge case is equality.
\end{proof}

Let $m\ge2$ and let $(P_1,\ldots,P_m)$ be any permutation of
$\{m+1,\ldots,2m\}$. Subscripts are cyclic, with $P_0=P_m$.
The order is $(1,P_1,2,P_2,\ldots,m,P_m)$. Define
\begin{equation}\label{eq:cells}
 a_i=\thR(P_{i-1},i),\quad b_i=\thR(i,P_i),\quad
 c_i=\thR(P_{i-1},P_i),\quad d_i=\max\{a_i+b_i,c_i\}.
\end{equation}
\begin{theorem}[Full Cell Criterion]\label{thm:cells}
The order is feasible exactly when $S_P(R):=\sum_i d_i\le2\pi$.
More precisely its feasible gap vectors are exactly
\begin{equation}\label{eq:gapset}
 x_i\ge a_i,\quad y_i\ge b_i,\quad x_i+y_i\ge c_i,
 \qquad\sum_i(x_i+y_i)=2\pi,
\end{equation}
where $x_i$ is the gap $P_{i-1}\to i$ and $y_i$ is $i\to P_i$.
The full minimum is the unique positive root $\rho_P$ of $S_P(R)=2\pi$.
\end{theorem}
\begin{proof}
Necessity follows from the two adjacent constraints and the high-pair
constraint on each whole cell. For sufficiency first take any positive
gaps satisfying the three local inequalities, without requiring closure.
For each pair choose either simple cyclic path; the following arguments
apply separately in both directions.

For two highs, the path consists of whole cells. Replacing their lengths
by $c_i$ and applying Lemma~\ref{lem:shell} on $[m+1,2m]$ gives the
pair bound. For a low $i$ and a high $P_j$, the first high on the path is
a neighbor $U$. If $U=P_j$, use adjacency. Otherwise the remaining
whole-cell high path has length at least $\thR(U,P_j)>\thR(i,P_j)$.
For two distinct lows $i,j$, let $U$ and $V$ be the first and last highs.
If $U=V$, the simple path has two edges and its first gap alone is at
least $\thR(i,U)>\thR(i,j)$. If $U\ne V$, the middle high path is at
least $\thR(U,V)>\thR(i,j)$. These cases exhaust all pairs and both arcs,
including paths crossing the cyclic cell $P_m\to1\to P_1$.
For $m=2$ each directed high path is one whole cell; no third high is
required. The low-high and low-low arguments remain valid.

When $S_P(R)\le2\pi$, take $x_i=a_i$, $y_i=d_i-a_i$ and add the
remaining $2\pi-S_P(R)$ to any single gap. Each directed path either
increases or stays fixed; both bounds survive and closure holds. This
also proves sufficiency of every vector in \eqref{eq:gapset}.
Each of $a_i+b_i,c_i$ is positive, continuous and strictly decreasing
in $R$; their maximum is strictly decreasing by choosing an attaining
branch at the larger radius. Moreover $S_P(0+)=2\pi m$ and
$S_P(\infty)=0$. The intermediate value theorem and monotonicity give
the unique root and its feasible placement.
\end{proof}

\section{Recovery of General Reflected Blocks}\label{sec:transfer}
Fix $0\le\alpha<1$ and nonnegative lengths $\ell_1,\ldots,\ell_k$,
$k\ge1$, with $T=\sum_j\ell_j<1-\alpha$. Put $t_0=0$,
$t_j=\sum_{h\le j}\ell_h$, $A=1+\alpha$, and
$h_\alpha(t)=1+\{t+\alpha\}$. Define a probability measure by
\begin{equation}\label{eq:mu}
\begin{split}
 \int F\dd\mu_{\alpha;\ell}
 ={}&\frac12\sum_{j=1}^k\int_{t_{j-1}}^{t_j}
 \bigl[F(t,A+t,A+t_{j-1}+t_j-t)\\
 &\hspace{38mm}+F(t,A+t_{j-1}+t_j-t,A+t)\bigr]\dd t\\
 &+\int_T^1F(t,h_\alpha(t),h_\alpha(t))\dd t.
\end{split}
\end{equation}
Each high marginal and the low marginal are uniform on their respective
intervals. Each slab reflection is a measure-preserving involution, and
the unchanged shift is a bijection modulo its wrap. Zero lengths and
endpoint assignments do not affect the measure. Adjacent slabs are
reflected separately, never as their union.

\begin{theorem}[Quantitative Full Transfer]\label{thm:transfer}
Let $I=\int g\dd\mu_{\alpha;\ell}$ and
$M_0=\max\{2,\lceil2/(1-\alpha-T)\rceil\}$.
There are genuine orders on $1,\ldots,2m$ with full minimum $\rho_m$
for every $m\ge M_0$, and, for $m\ge\max\{M_0,2048\}$,
\begin{equation}\label{eq:radiuserror}
 \left|\frac{\rho_m}{(2m)^2}-\frac I{4\pi}\right|
 \le\frac1{4\pi}\left(\frac{6k^2+28k+1060}{m}
                         +\frac{16384}{3m^2}\right).
\end{equation}
Consequently $\limsup_n\Rstar(n)/n^2\le C(\mu_{\alpha;\ell})$.
No chord-dominance assumption on any slab is required.
\end{theorem}
\begin{proof}
Define exact integers
\[
 s=\lfloor\alpha m\rfloor,\quad l_j=2\lfloor\ell_jm/2\rfloor,
 \quad a_0=0,\quad a_j=\sum_{h\le j}l_h,\quad z=a_k,
 \quad r=m-s,
\]
and $H(u)=m+1+((u+s-1)\bmod m)$. On the block $(a_{j-1},a_j]$ put
$J(i)=2a_{j-1}+l_j+2-i$ when $i$ is even, and $J(i)=i$ otherwise.
Set $P_i=H(J(i))$, $P_0=P_m$. Each block reverses its even ranks and
fixes its odd ranks, giving a disjoint involution; $H$ is a bijection.
Lengths zero and two cause no exception. Since
$s+z\le m(\alpha+T)$, $r-z\ge2$. The blocks and their exits precede
the high wrap.

For completeness the cells have the following exhaustive inventory.
A nonempty block $(a,a+l]$ has interior indices $a+2,\ldots,a+l$.
The predecessor/current pair there is
\[
 \begin{cases}
 (m+s+i-1,m+s+2a+l+2-i),&i\text{ even},\\
 (m+s+2a+l+3-i,m+s+i),&i\text{ odd}.
 \end{cases}
\]
The exceptional set is
$\{1,r,r+1\}\cup\{a_j+1:l_j>0\}$, intersected with
$\{1,\ldots,m\}$, with duplicates counted once. It has at most $k+3$
members. Cell $1$ is $(H(m),H(1))$; the exit from $(a,a+l]$ is
$(H(a+2),H(a+l+1))$. It is also the next nonempty block's entry even
if zero lengths intervene. Cell $r$ is $(2m-1,2m)$, and cell $r+1$,
when present, is $(2m,m+1)$. All other cells are ordinary
$(H(i-1),H(i))$. The interiors exclude every shared entry index, and
the wrap lies outside them, proving disjointness and completeness.

Let $\mu_m=m^{-1}\sum_i\delta_{(i/m,P_{i-1}/m,P_i/m)}$.
For a continuous, possibly nonsymmetric $F$, write
$B_F=\norm F_\infty$ and $\omega_F$ for its max-norm modulus.
Replace $\alpha,\ell_j$ in \eqref{eq:mu} by $s/m,l_j/m$ to obtain
$\bar\mu_m$. On each double panel
$[(a+2h-2)/m,(a+2h)/m]$, assign half the first-orientation integral
to its even cell and half the reverse-orientation integral to its odd
cell. Each mass is $1/m$. Interior coordinate errors are at most $4/m$;
ordinary diagonal panels have errors at most $1/m$. Exceptional panel
assignments each cost at most $2B_F/m$, so
\[
 \left|\int F\dd(\mu_m-\bar\mu_m)\right|
 \le\omega_F(4/m)+2(k+3)B_F/m.
\]
The floor errors are less than $1/m$ for the shift and $2/m$ for each
length. The intervals between $a_j/m$ and $t_j$, together with the wrap
interval, have total length at most $[k(k+1)+1]/m$. This bound allows
overlap and zero slabs. Off these intervals the same slab and orientation
apply; the reflected coordinate error in slab $j$ is less than
$(4j-1)/m$ and the other is less than $1/m$. Thus
\begin{equation}\label{eq:weakerror}
 \left|\int F\dd(\mu_m-\mu)\right|
 \le\omega_F(4/m)+\omega_F((4k-1)/m)
       +\frac{(2k^2+4k+8)B_F}{m}.
\end{equation}

The complete maximum $g$ obeys $1\le g<3$ and is $4$-Lipschitz.
For $t\le1/4$ the chord dominates because
$\sqrt{t/X}+\sqrt{t/Y}\le2\sqrt t\le1$. Above $1/4$, the sum of
absolute partial derivatives of the chain is at most $2\sqrt2+1<4$,
and that of the chord is less than two. The maximum of these Lipschitz
functions has the same bound; splitting at $t=1/4$ proves it globally.
Putting $G_m=\int g\dd\mu_m$, \eqref{eq:weakerror} yields
\begin{equation}\label{eq:scoreerror}
 |G_m-I|\le(6k^2+28k+36)/m.
\end{equation}
This estimate retains every finite maximum and seam even if its branch
differs from its continuum counterpart.

By Theorem~\ref{thm:cells}, the actual full minimum is the root of
$S_m(R)=\sum_i\max\{\thR(P_{i-1},i)+\thR(i,P_i),
\thR(P_{i-1},P_i)\}=2\pi$. For $R=4cm^2$, $c\ge1/32$, $m\ge32$,
and $1\le u,v\le2m$, put $q=\sqrt{uv}/R$. Then $q\le1/(2cm)\le1/2$.
The relative denominator change satisfies
\[
 0\le1-\frac1{\sqrt{(1+u/R)(1+v/R)}}\le\frac{u+v}{2R}.
\]
Also $0\le\arcsin t-t\le t^3/3$ for $0\le t\le1/2$,
by integrating $1/\sqrt{1-t^2}-1\le t^2$ on that interval.
It follows that
\[
 \left|\thR(u,v)-\frac{2\sqrt{uv}}R\right|
 \le\frac{512}{m^2}+\frac{8192}{3m^3},\qquad
 \left|S_m(4cm^2)-\frac{G_m}{2c}\right|
 \le E_m:=\frac{1024}m+\frac{16384}{3m^2}.
\]
The last step uses the one-Lipschitz property of the scalar maximum and
two angles per chain. For $m\ge2048$, $E_m<1$; at $c=1/32$ the score
exceeds $15>2\pi$, while at $c=1/2$ it is below $4<2\pi$.
Thus the actual normalized root $c_m=\rho_m/(2m)^2$ is first bracketed
in $(1/32,1/2)$. Substitution then gives
$|4\pi c_m-G_m|\le2c_m E_m\le E_m$.
Together with \eqref{eq:scoreerror} this is \eqref{eq:radiuserror}.
The root placement is all-pairs feasible. Deleting only label $2m$
keeps exactly $1,\ldots,2m-1$ and all their constraints; the normalization
factor $(2m/(2m-1))^2$ tends to one. This proves the global limsup
bound for both parities without identifying an odd fixed-order minimum.
\end{proof}

The same construction allows countably many lengths with total
$T<1-\alpha$. Truncation after $k$ blocks and reversion of the tail to
the diagonal changes any $F$ integral by at most
$2B_F\sum_{j>k}\ell_j$. For $g$ the bound is six times that tail length.
Apply the finite theorem and then let $k\to\infty$. A single recovery
sequence also results by choosing increasing thresholds $M_k$ growing
faster than $k^3$, each satisfying the common geometric gate, and using the largest $k\le m$ with $M_k\le m$.
Both \eqref{eq:weakerror} for Lipschitz tests and the omitted tail vanish;
density of Lipschitz functions on compact $\mathcal D$ gives continuous-test
recovery. This supplies countable upper constructions, not an optimum
over partitions.

\section{The Exact Normalized Prefix Parameter}\label{sec:x}
For $0\le u\le x\le1$, define
\begin{align*}
 c(u,x)&=\sqrt{(1+u)(1+x-u)},\\
 k(u,x)&=\sqrt u(\sqrt{1+u}+\sqrt{1+x-u}),\\
 d(u)&=\max\{1+u,2\sqrt{u(1+u)}\}.
\end{align*}
\begin{equation}\label{eq:E}
 E(x)=\int_0^x[\max\{c(u,x),k(u,x)\}-d(u)]\dd u.
\end{equation}
\begin{theorem}\label{thm:x}
The function $E$ has a unique global minimum $x_*$ on $[0,1]$, and
$719/2500<x_*<2877/10000$. These are the exact bounds
$0.2876<x_*<0.2877$.
\end{theorem}
\begin{proof}
The ratio $v(u,x)=k/c=\sqrt{u/(1+u)}+\sqrt{u/(1+x-u)}$
strictly increases in $u$. Its endpoint equals one at the unique
$\tau\in(1/4,1/3)$ solving
$\sqrt{\tau/(1+\tau)}+\sqrt\tau=1$.
Set $z(x)=x$ for $x\le\tau$; for $x>\tau$ let $z\in(0,x)$ be
the unique unsquared root $v(z,x)=1$. The block is chord up to $z$,
then chain. The diagonal switches at $u=1/3$.
For $x>\tau$, with $y=1+x-z$ and $r=\sqrt{z/(1+z)}$,
\begin{equation}\label{eq:zprime}
 z'=\frac{z}{(y/(1+z))^{3/2}+1+x}
    =\frac{(1-r)^2}{2-r+2r^2}>0.
\end{equation}
The second equality uses $y=z/(1-r)^2$ and $z=r^2/(1-r^2)$.
All denominators are positive.

Let $\Phi=E'$. Cost equality at the switch cancels its first moving
boundary term. The endpoint ties at $x=\tau,1/3$ also agree, so $E$
is $C^1$ there. On the chord interval,
\begin{gather*}
 \Phi(x)=\sqrt{1+x}+(1+x/2)\arcsin\frac{x}{2+x}-(1+x),\\
 \Phi(0)=\Phi'(0)=0,\qquad
 \Phi''(x)=-\frac1{2(2+x)(1+x)^{3/2}}<0.
\end{gather*}
Thus $\Phi<0$ on $(0,\tau]$. On $x>\tau$, put
\[
 I_1=\int_0^z\sqrt{\frac{1+u}{1+x-u}}\dd u,\qquad
 I_2=\int_z^x\sqrt{\frac u{1+x-u}}\dd u,\qquad
 p=\sqrt x(\sqrt{1+x}+1).
\]
Differentiation yields
\begin{equation}\label{eq:Phi}
 \Phi=p-d(x)+(I_1+I_2)/2,\qquad
 \Phi'=p'-d'(x)+\sqrt x/2-J/4+Q,
\end{equation}
where
\[
 J=\int_0^z\frac{\sqrt{1+u}}{(1+x-u)^{3/2}}\dd u
   +\int_z^x\frac{\sqrt u}{(1+x-u)^{3/2}}\dd u,\quad
 Q=\frac{z'(\sqrt{1+z}-\sqrt z)}{2\sqrt{1+x-z}}>0.
\]
For $\tau<x<1/3$, $d'=1$, $J\le x\sqrt{1+x}<7/18$ and
$p'>1/\sqrt x\ge\sqrt3>5/3$. Hence $\Phi'>41/72>0$.
At $x=1/3$, $v(2/7,x)<1$: the pre-square residual is $50/99>0$
and its square minus $4(2/9)(3/11)$ is $124/9801>0$.
Thus $z>2/7$, and
\[
 \Phi(1/3)>\frac1{\sqrt3}-\frac23+\frac{\sqrt3}{14}
 =\frac{17\sqrt3}{42}-\frac23>0
\]
by $17^2\cdot3>28^2$. There is exactly one zero $x_*$ in
$(\tau,1/3)$, with strict decrease then increase of $E$ on $[0,1/3]$.

Global minimality needs a separate tail argument; convexity must not be
extrapolated across $1/3$. On $1/3<x\le1$, \eqref{eq:Phi} becomes
\[
 \Phi'=-M-J/4+Q,\qquad
 M=\frac{1+2x-(1+x)^{3/2}}{2\sqrt{x(1+x)}}.
\]
Since $z>2/7$, $r>\sqrt{2/9}>8/17$. The polynomial
$5-13r+5r^2$ decreases on $(0,1)$ and is $-3/289$ at $8/17$;
\eqref{eq:zprime} gives $z'<1/7$. Also
$\sqrt{1+z}-\sqrt z<2/3$, so $Q<1/21$.
Write $q=\sqrt{1+x}\in[2/\sqrt3,\sqrt2]$. Then
\[
 M=\sqrt{\frac{q-1}{q+1}}\frac{1+q-q^2}{2q}.
\]
For $q\le4/3$ the factors exceed $1/4$ and are at least $5/24$,
respectively, so $M>5/96>1/20$. For $q\ge4/3$ they are at least
$1/\sqrt7>3/8$ and $(2-\sqrt2)/4>1/7$, so again $M>1/20$.
These comparisons follow by positive squares, using
$2/\sqrt3>17/15$ and $\sqrt2<10/7$. Therefore
$\Phi'<-1/420$ throughout the tail.
The rational gates in Section~\ref{sec:gates} give $E(1)>3/250>0$.
Strict concavity makes $E$ at least its endpoint chord on $[1/3,1]$.
Both endpoint values exceed $E(x_*)<0$, proving unique global
minimality. The same gates give opposite signs at $719/2500$ and
$2877/10000$, both above $\tau$, locating this unique root. They also
give $\Phi(4/5)<0$; thus a later local maximum exists, and does not
invalidate the endpoint-chord argument.
\end{proof}

\section{The Exact Shift Parameter}\label{sec:alpha}
For $0\le\alpha\le1/2$ let $A=1+\alpha$, $a=A/3$, $b=1-\alpha$,
and
\[
 K(\alpha)=\int_0^1\max\{\sqrt{t h_\alpha(t)},h_\alpha(t)/2\}\dd t,
 \qquad e=E(x_*),\qquad F(\alpha)=K(\alpha)+A^2e/2.
\]
The full prefix cost at length $Ax_*$ is $F(\alpha)/(2\pi)$.
Indeed $Ax_*<a\le b$, the replaced diagonal is chord, and scaling
$t=Au$ gives its increment $A^2E(x_*)$. Moreover
$b-Ax_*>1/2-(3/2)(2877/10000)=1369/20000>1/15$,
so the reflection stays before the wrap even at the closed endpoints.

\begin{theorem}\label{thm:alpha}
There is exactly one zero $\al$ of
$F'(\alpha)=K'(\alpha)+(1+\alpha)e$ on $[0,1/2]$.
It is the unique minimum of this fixed-$x_*$ prefix family, and
$1093/10000<\al<10931/100000$.
\end{theorem}
\begin{proof}
The split integral is
\[
 K=\int_0^a(A+t)/2\dd t+\int_a^b\sqrt{t(A+t)}\dd t
                       +\int_b^1\sqrt{t(\alpha+t)}\dd t.
\]
Differentiation, including the moving wrap, gives
\begin{equation}\label{eq:D}
 D(\alpha):=K'(\alpha)
 =\frac a2+\frac12\int_a^b\sqrt{\frac t{t+A}}\dd t
 +\frac12\int_b^1\sqrt{\frac t{t+\alpha}}\dd t
 -(\sqrt2-1)\sqrt b.
\end{equation}
For $0<x<1/3$, centering the chord at $M=1+x/2$ gives
$E(x)\ge-\int_{-x/2}^{x/2}u^2/M\dd u=-x^3/(12M)>-1/324$.
At $x=1/4$, which is strictly chord, rationalization gives
$-E(1/4)>1/1728$. Theorem~\ref{thm:x} implies
$-1/324<e<-1/1728$.
A second differentiation of \eqref{eq:D} gives
\[
 K''=\frac1{12}+\frac{(1-1/\sqrt2)\sqrt b}{2}
       +\frac{\sqrt2-1}{2\sqrt b}-\frac{T_1}{4},
 \quad
 T_1=\int_a^b\frac{\sqrt t}{(t+A)^{3/2}}\dd t
       +\int_b^1\frac{\sqrt t}{(t+\alpha)^{3/2}}\dd t.
\]
Both denominator bases are at least one; $T_1\le1-a\le2/3$.
Using $1/2\le b\le1$ and $\sqrt2>7/5$ yields $K''>7/60$,
so $F''>7/60-1/324=46/405>1/9$ on the whole domain.

At zero,
$D(0)=5/6-\sqrt2/2-\tfrac12\log((1+\sqrt2)/\sqrt3)$.
For $z>1$, $\log z>2(z-1)/(z+1)$ by differentiation. Since the
argument exceeds $48/35$, using $\sqrt3<7/4$ gives
$F'(0)<5/6-7/10-13/83=-29/1245<0$.
At $1/2$ the surviving integrand in \eqref{eq:D} is at least
$1/\sqrt2$, so $D(1/2)\ge5/(4\sqrt2)-3/4>1/8$ and
$F'(1/2)>1/8-(3/2)/324=13/108>0$.
Strict convexity proves uniqueness and minimality.

For the tighter location set $x_0=5753/20000$, the midpoint of the
$x_*$ bracket. On $(\tau,1/3)$, \eqref{eq:Phi} gives
$E''<8/3-1+1/3+1/8=17/8<3$: here $p'<8/3$,
$z'<x/(1+x)<1/4$, $Q<1/8$, and $J\ge0$.
Taylor's theorem at the stationary minimum gives
\begin{equation}\label{eq:ebox}
 E(x_0)-3/800000000<e\le E(x_0).
\end{equation}
The directed gates in Section~\ref{sec:gates}, evaluating only this
rational argument and the two rational $\alpha$ endpoints, give
$F'(1093/10000)<0<F'(10931/100000)$. These isolate the same unique zero.
\end{proof}

\section{A Common Width Lemma and the Next Two Parameters}\label{sec:widths}
Fix $A>1$ and $A/4<u<A/3$. Put $a=A/3$, $B=A+u$, and $h=a-u$,
so $B>5A/4$, $0<h<A/12$, $h/B<1/15$. Reflect a diagonal slab
$[u,u+d]$, $0\le d\le h$, separately from all preceding slabs. Its
unnormalized complete-cost increment is exactly
\begin{equation}\label{eq:widthcost}
\begin{split}
 D_{A,u}(d)=\int_0^d\bigl[\max\{&\sqrt{(B+s)(B+d-s)},\\
 &\sqrt{u+s}(\sqrt{B+s}+\sqrt{B+d-s})\}-(B+s)\bigr]\dd s.
\end{split}
\end{equation}
The removed diagonal is chord because $A-3(u+s)\ge0$; its only possible
tie is $s=d=h$. Reflection preserves each marginal as in
\eqref{eq:mu}, and unchanged slabs cancel with their full maxima.

\begin{lemma}[Unique Mixed Width]\label{lem:width}
The function $D_{A,u}$ has a unique minimum $d_*$ on $[0,h]$.
Define $\tau_u$ as the unique unsquared root in $(0,h)$ of
\[
 W(d)=\sqrt{\frac{u+d}{B+d}}+\sqrt{\frac{u+d}B}=1.
\]
Then $\tau_u<d_*<\tau_u+\tau_u^2/(8B)$, and $d_*$ is the unique
mixed-branch stationary point. The cost decreases strictly up to $d_*$
and increases strictly afterwards.
\end{lemma}
\begin{proof}
Set $t=u+s$, $X=B+s$, $Y=B+d-s$, $c=\sqrt{XY}$,
$k=\sqrt t(\sqrt X+\sqrt Y)$, and $V=k/c$.
Then
\[
 V_s=\frac{A/X^{3/2}+(Y+t)/Y^{3/2}}{2\sqrt t}>0.
\]
For $d>0$, $V(0,d)<1$ and
$V(d/2,d)=2\sqrt{(u+d/2)/(A+u+d/2)}<1$.
The endpoint $W$ strictly increases, with $W(0)<1$ and
$W(h)=1/2+\sqrt{a/B}>1$. This proves the asserted entry width.
For $d\le\tau_u$ set $z=d$. For $d>\tau_u$ there is a unique
$z\in(d/2,d)$ with $V(z,d)=1$. The cost is
\[
 D_{A,u}=\int_0^z c\dd s+\int_z^d k\dd s-Bd-d^2/2.
\]
Let $\Psi=D_{A,u}'$. In the all-chord regime,
\[
 \Psi(d)=B\phi(d/B),\qquad
 \phi(q)=\sqrt{1+q}+(1+q/2)\arcsin\frac q{2+q}-(1+q).
\]
Here $\phi(0)=\phi'(0)=0$ and
$-1/4<\phi''(q)=-[2(2+q)(1+q)^{3/2}]^{-1}<0$ for $q>0$.
Therefore $-\tau_u^2/(8B)<\Psi(\tau_u)<0$ and the positive-width
chord cost is strictly decreasing.

In the mixed regime put $q=u+d$,
$p=\sqrt q(\sqrt{B+d}+\sqrt B)$,
$I_c=\int_0^z\sqrt{X/Y}\dd s$ and $I_k=\int_z^d\sqrt{t/Y}\dd s$.
Leibniz differentiation gives
\begin{equation}\label{eq:Psi}
 \Psi=p-(B+d)+(I_c+I_k)/2,\qquad
 z'=\frac{t}{A(Y/X)^{3/2}+Y+t}\in(0,1/2),
\end{equation}
where $t,X,Y$ in $z'$ are evaluated at $s=z$.
The first moving-boundary term cancels by $c(z)=k(z)$.
The bound on $z'$ follows from $Y-t=A+d-2z>A-d>0$.
The cost and first derivative are continuous at entry: the chain interval
shrinks to zero and the endpoint costs agree. Equivalently put $s=dr$
on $[0,1]$; positive radicands bound the difference quotients and the tie
set has measure zero, permitting dominated differentiation. This also
justifies the one-sided derivative at $h$.

The second derivative retains a nonzero moving-switch term:
\begin{equation}\label{eq:Psiprime}
 \Psi'=p'-1+\tfrac12\sqrt{q/B}-J/4+Q,
 \quad p'=\frac{A+2q}{2\sqrt{q(A+q)}}+\frac{\sqrt B}{2\sqrt q},
\end{equation}
where
\[
 J=\int_0^z\frac{\sqrt X}{Y^{3/2}}\dd s+
       \int_z^d\frac{\sqrt t}{Y^{3/2}}\dd s,
 \qquad Q=\frac{z'(\sqrt{B+z}-\sqrt{u+z})}{2\sqrt{B+d-z}}>0.
\]
The first term in $p'$ decreases in $q$ and equals $5/4$ at $q=a$.
Thus $p'>5/4+\sqrt{15}/4>17/8$.
Also $0<J\le d\sqrt{B+d}/B^{3/2}<2/15$, giving
$\Psi'>131/120>1$ everywhere in the mixed regime.
Cost equality cannot be used to drop $Q$, since the derivative integrands
at the switch do not coincide.

At $h$, $B=4a-h$, $u=a-h>3a/4$, and
\[
 \Psi(h)=\sqrt{aB}-2a+(I_c+I_k)/2,\qquad
 2a-\sqrt{aB}=\frac h{2+\sqrt{B/a}}<2h/7.
\]
On this slab $15a/4<X,Y\le4a$ and $3a/4<t\le a$, so
$\sqrt{X/Y}>9/10$ and $\sqrt{t/Y}>2/5$.
Since $z(h)>h/2$, $(I_c+I_k)/2>13h/40$, whence
$\Psi(h)>11h/280>0$. Continuity and strict increase of $\Psi$ in the
mixed regime give its unique zero $d_*$. The chord sign excludes earlier
stationary points, proving the unique global minimum including both
endpoints. Finally integrating $\Psi'>1$ from $\tau_u$ to $d_*$ gives
$d_*-\tau_u<-\Psi(\tau_u)<\tau_u^2/(8B)$.
\end{proof}

We now fix $A=1+\al$ and $\lambda=Ax_*$. Define
\begin{equation}\label{eq:epsdefinition}
 \epsb=\mathop{\rm argmin}_{0\le d\le A/3-\lambda}D_{A,\lambda}(d).
\end{equation}
Lemma~\ref{lem:width} proves this is a unique mixed stationary minimum.
The rational gates below give
\[
 43/1000<\tau_\lambda<87/2000,
 \quad 43/1000<\epsb<87/2000+(87/2000)^2/8<7/160.
\]
In particular $v=\lambda+\epsb$ remains in $(A/4,A/3)$.
Define in turn
\begin{equation}\label{eq:Deltadefinition}
 \Ds=\mathop{\rm argmin}_{0\le d\le A/3-v}D_{A,v}(d).
\end{equation}
The same lemma proves its uniqueness and identifies its stationarity
with $4\pi C_3'(\Ds)=0$. The second pair of rational gates gives
\begin{equation}\label{eq:Deltabracket}
\begin{gathered}
 29/5000<\tau_v<67/10000,\\
 29/5000<\Ds<67/10000+(67/10000)^2/8<27/4000<A/3-v.
\end{gathered}
\end{equation}
These definitions retain all mixed branches. They do not use a chord-only
stationary equation or a rounded parameter.

\section{Rational Gates and Their Reproduction}\label{sec:gates}
This section supplies the finite arithmetic premises used above. Every
checker in the accompanying \texttt{checks} directory is self-contained;
none imports production code, the independent geometric verifier, another
checker, or task records. The first two have an optional numerical mode;
use \texttt{--exact-only} for the theorem gates. Run with ordinary Python
assertions enabled (do not use \texttt{-O}).

\subsection{Prefix Gates}
The file \texttt{upper\_lambda.py} uses interval endpoints on the grid
$10^{-50}$. Each arithmetic operation rounds outwards, division across
zero is rejected, and $\sqrt{n/10^{50}}$ is enclosed using the integer
square root of $n10^{50}$. Two half-angle reductions express
$\arctan y=4\arctan q$ with $|q|<1$. The first 64 alternating terms
are widened by $|q|^{129}/129$ before multiplying by four.
For $\log y$, take two square roots, set
$q=(y^{1/4}-1)/(y^{1/4}+1)$, and use the 64-term odd series for
$\log y=8\sum_{j\ge0}q^{2j+1}/(2j+1)$, with the corresponding
tail bounded by $8|q|^{129}/[129(1-|q|^2)]$.
Every operation in those formulas is itself enclosed.

At rational $x>\tau$, 100 dyadic bisections of $[0,x]$ isolate the
unsquared equation $v(z,x)=1$. An unseparated sign aborts the check.
The exact primitives
\[
 \mathcal F(t,c)=\frac{(2t+c)\sqrt{t(t+c)}
     -c^2\log((\sqrt t+\sqrt{t+c})/\sqrt c)}4,
\]
\[
 \mathcal J(v,M)=\frac{v\sqrt{M^2-v^2}
                +M^2\arctan(v/\sqrt{M^2-v^2})}2,
 \qquad
 \mathcal H(c,t)=c\arctan\sqrt{t/(c-t)}-\sqrt{t(c-t)}
\]
have derivatives $\sqrt{t(t+c)}$, $\sqrt{M^2-v^2}$ and
$\sqrt{t/(c-t)}$, respectively. They evaluate the split integrals in
\eqref{eq:E} and \eqref{eq:Phi}. Explicitly the block integral is
\begin{align*}
 &\mathcal J(z-x/2,1+x/2)-\mathcal J(-x/2,1+x/2)\\
 &\quad+\mathcal F(x,1)-\mathcal F(z,1)\\
 &\quad+\mathcal J(x-(1+x)/2,(1+x)/2)
 -\mathcal J(z-(1+x)/2,(1+x)/2).
\end{align*}
The last four terms are omitted in the chord regime. Subtract
$x+x^2/2$ for $x\le1/3$, or
$7/18+2[\mathcal F(x,1)-\mathcal F(1/3,1)]$ otherwise.
For \eqref{eq:Phi},
$I_1=\mathcal H(2+x,1+z)-\mathcal H(2+x,1)$ and
$I_2=\mathcal H(1+x,x)-\mathcal H(1+x,z)$.
This gives the strict rational gates
\begin{center}
\begin{tabular}{lll}\toprule
Quantity & Lower Bound & Upper Bound\\\midrule
$\Phi(719/2500)$ & $-1/20000$ & $-1/25000$\\
$\Phi(2877/10000)$ & $1/10000$ & $11/100000$\\
$\Phi(4/5)$ & $-1/500$ & $-1/1000$\\
$E(1)$ & $3/250$ & $13/1000$\\\bottomrule
\end{tabular}
\end{center}
Direct directed square roots also show $719/2500>\tau$.
These four gates establish constant signs, not the continuum derivative
inequalities proved in Section~\ref{sec:x}.

\subsection{Shift Gates}
The file \texttt{upper\_alpha\_minimum.py} uses exact rational arithmetic
and directed roots on the grid $10^{-24}$, checking each endpoint by
squaring. It isolates $z(x_0)$ with 64 sign-separated bisections.
On $[0,z_l]$ it integrates $c$ and on $[z_h,x_0]$ it integrates $k$,
using 2048 panels each. Both integrands are concave: their second
derivatives are $-M^2/c^3$ for the chord and the sum of
$-1/[4(u(1+u))^{3/2}]$ and
$-(1+x_0)^2/[4(u(1+x_0-u))^{3/2}]$ for the chain.
Thus composite trapezoids are lower bounds and midpoints are upper bounds.
The omitted interval contributes between zero and $4(z_h-z_l)$.
Subtract the exact diagonal and apply \eqref{eq:ebox}.
For each $D$ integral in \eqref{eq:D}, use 256 panels and the concavity
\[
 \frac{\mathrm d^2}{\mathrm dt^2}\sqrt{\frac t{t+c}}
 =-\frac{c(4t+c)}{4t^{3/2}(t+c)^{5/2}}<0.
\]
The negative wrap term is bounded in the reversed direction.
The resulting outward intervals have common denominator $10^{12}$:
\begin{center}
\begin{tabular}{lrr}\toprule
Quantity & Lower Numerator & Upper Numerator\\\midrule
$E(x_0)$ & $-844268665$ & $-844268070$\\
$e$ & $-844272415$ & $-844268070$\\
$D(1093/10000)$ & $935124205$ & $935201790$\\
$F'(1093/10000)$ & $-1427184$ & $-1344781$\\
$D(10931/100000)$ & $938842182$ & $938919761$\\
$F'(10931/100000)$ & $2282349$ & $2364747$\\\bottomrule
\end{tabular}
\end{center}
No evaluation at an uncertified irrational root or floating quadrature
is used. The analytic convexity proof supplies uniqueness.

\subsection{Second and Third Width Gates}
Write $\alpha_L=1093/10000$, $\alpha_H=10931/100000$,
$x_L=719/2500$, $x_H=2877/10000$, and $A_L=1+\alpha_L$,
$A_H=1+\alpha_H$.
For the second width, put
$f(x,q)=\sqrt{(x+q)/(1+x+q)}+\sqrt{(x+q)/(1+x)}$.
It increases in $q$ and, for $q<1$, in $x$; the squared summands'
$x$ derivatives are $1/(1+x+q)^2$ and $(1-q)/(1+x)^2$.
Consequently
\[
 W(43/1000)<f(x_H,(43/1000)/A_L)<1,
 \quad W(87/2000)>f(x_L,(87/2000)/A_H)>1.
\]
For positive rational squared summands $p,q$ with residual $r=1-p-q>0$,
\begin{equation}\label{eq:squaregate}
 \operatorname{sign}(\sqrt p+\sqrt q-1)
 =\operatorname{sign}(4pq-r^2).
\end{equation}
The file \texttt{upper\_boundary.py} checks the positive residuals and
the respective exact signed differences
\[
 -\frac{342531413518344092144807428477}{441789962670730640526171361763841},
 \qquad
 \frac{51974394530491111127115563}{253438432681061908372345318521}.
\]
Also $B>1$ and
$7/160-87/2000-(87/2000)^2/8=431/32000000>0$,
which proves the sharpened upper bound for $\epsb$.

Put $e_L=43/1000$, $e_U=7/160$ and
\[
 V_L=A_Lx_L+e_L=9050867/25000000,\qquad
 V_U=A_Hx_H+e_U=362898487/10^9.
\]
Thus $V_L<v<V_U$. For
$W(A,v,d)=\sqrt{(v+d)/(A+v+d)}+\sqrt{(v+d)/(A+v)}$,
both squared summands decrease in $A$ and increase in $v$ when $d<A$.
The file \texttt{upper\_mixed\_width.py} checks
\[
 W(A_L,V_U,29/5000)<1,\qquad W(A_H,V_L,67/10000)>1.
\]
Their signed differences \eqref{eq:squaregate}, with positive residuals,
are respectively
\[
 -\frac{963745517404602389188679409068317}{4734571866017504812540302482915910561},
 \qquad
 \frac{34763358497298510706339907}{205268116362886684743388903969}.
\]
Finally
\[
 \frac{27}{4000}-\frac{67}{10000}-\frac{(67/10000)^2}{8}
 =\frac{35511}{800000000}>0,\qquad
 \frac{A_L}3-V_U-\frac{27}{4000}
 =\frac{354539}{3000000000}>0.
\]
These show that all corner evaluations are admissible and prove
\eqref{eq:Deltabracket}. The same margins give $A/4<v<A/3$ as required
when Lemma~\ref{lem:width} is applied a second time.

\section{The Fourth Block and the Final Upper Endpoint}\label{sec:fourth}
Fix all four implicit parameters above and write $B_4=A+w$,
$\eta_0=1/20000$, $M=B_4+\eta/2$.
The independent rational enclosing corner
$W_U=V_U+27/4000=369648487/10^9$ gives, for $0<\eta\le\eta_0$,
\begin{align}
 A-3w&>354539/10^9, & A-3w-4\eta&>154539/10^9,\nonumber\\
 A/3-w-\eta&>204539/(3\cdot10^9),
 &1-\al-w-\eta&>520991513/10^9,\label{eq:fourthgates}\\
 3/2-M&>21016513/10^9.&&\nonumber
\end{align}
Thus the fourth slab lies strictly before the diagonal cutoff and wrap.
Only this slab changes. For every continuous, possibly nonsymmetric $F$,
its measure increment is
\begin{multline*}
 \frac12\int_w^{w+\eta}
 [F(t,A+t,A+2w+\eta-t)+F(t,A+2w+\eta-t,A+t)\\
 -2F(t,A+t,A+t)]\dd t.
\end{multline*}
This preserves positivity and each marginal. The shared endpoint has zero
mass; the third and fourth reflections remain separate. No stationarity
of earlier blocks is used to cancel their unchanged full costs.

At $t=w+s$, $X=B_4+s$, $Y=B_4+\eta-s$, $0\le s\le\eta$,
\eqref{eq:fourthgates} implies
\[
 \frac{\sqrt t(\sqrt X+\sqrt Y)}{\sqrt{XY}}
 \le2\sqrt{(w+\eta)/B_4}<1,
 \qquad X-4t=A-3w-3s>0.
\]
Both the inserted and removed full maxima are strictly chord throughout
the closed slab. The unnormalized increment is therefore
\begin{align}
 4\pi[C_4(\eta)-C_3(\Ds)]
 &=\int_0^\eta[\sqrt{(B_4+s)(B_4+\eta-s)}-(B_4+s)]\dd s\nonumber\\
 &=-\int_{-\eta/2}^{\eta/2}
       \frac{z^2}{M+\sqrt{M^2-z^2}}\dd z.\label{eq:decrement}
\end{align}
The radicand is at least $B_4(B_4+\eta)>0$. Equivalently this is
\[ (\eta/2)\sqrt{B_4(B_4+\eta)}+M^2\arcsin(\eta/(2B_4+\eta))-\eta M. \]
Since
\[
 2B_4\le M+\sqrt{M^2-z^2}\le2M,\qquad
 \int_{-\eta/2}^{\eta/2}z^2\dd z=\eta^3/12,
\]
we obtain, with $J_4=C_4-C_3(\Ds)$,
\[
 -\frac{\eta^3}{96\pi B_4}\le J_4(\eta)
 \le-\frac{\eta^3}{96\pi M}< -\frac{\eta^3}{576}<0,
\]
using $M<3/2$ and $\pi<4$. Furthermore
\[
 J_4(\eta)=-\frac{\eta^3}{96\pi B_4}+\mathcal R_4(\eta),
 \qquad0\le\mathcal R_4(\eta)\le\frac{\eta^4}{192\pi B_4^2}.
\]
Indeed the remainder is at most $\eta^3(1/B_4-1/M)/(96\pi)\le\eta^4/(192\pi B_4^2)$. This is a signed finite-interval remainder and proves the first nonzero
unilateral term is cubic. It is not a negative first variation.
At $\eta_0$, the strict saving is
$\eta_0^3/576=1/4608000000000000$.

The four lengths sum to $T=w+\eta_0<1-\al$, so
Theorem~\ref{thm:transfer} applies with $k=4$. Its radius error is
$[1268/m+16384/(3m^2)]/(4\pi)$. The actual full roots and the odd
placements obtained by deletion prove \eqref{eq:target}. This establishes
the upper construction independently of any finite-size optimization.
It does not establish optimality of four blocks, of the marginal
relaxation, or of the global geometric problem.

\subsection{Independent Finite Check of the Fourth Saving}
The file \texttt{upper\_fourth\_block.py} verifies
\eqref{eq:fourthgates} with fractions and uses a 16-panel midpoint upper
enclosure of the raw inserted full maximum, subtracting the exact
diagonal integral. The branch gates above justify concavity in $s$.
For the chord difference its derivative in $B_4$ is
\[
 \frac{2B_4+\eta}{2\sqrt{(B_4+s)(B_4+\eta-s)}}-1\ge0,
\]
since the numerator's square minus four times the radicand is
$(\eta-2s)^2$. Therefore the upper enclosing $B_4$ corner gives an
upper cost bound. Directed roots on a $10^{-40}$ grid make this an exact
rational integral enclosure, independent of the antiderivative and the
rationalized form in \eqref{eq:decrement}. The code also checks 16 corners,
448 polynomial reflection moments and 48 branch probes, as bounded
bookkeeping diagnostics. It rejects zero, negative, and cutoff widths
for the strict chord-slab gate. Those probes are not the proofs of
arbitrary-test marginal preservation, universal recovery, or the cubic
coefficient; the preceding analytic arguments supply those quantifiers.

\section{Scope, Provenance, and Reproduction}
This text repackages existing exact theorems and proved corollaries from
the seven research sources indexed in \texttt{UPPER\_DEPENDENCIES.md}.
That document maps each section and checker to its original source at
the frozen scientific snapshot. The two width arguments are written once
with an explicit generic start, and checked under the hypotheses of both
applications. No new mathematical coefficient, finite certificate, or
claim of external mathematical acceptance is introduced.

From the directory containing this source, run
\begin{verbatim}
python -I -S checks/upper_lambda.py --exact-only
python -I -S checks/upper_alpha_minimum.py --exact-only
python -I -S checks/upper_boundary.py
python -I -S checks/upper_mixed_width.py
python -I -S checks/upper_fourth_block.py
\end{verbatim}
The first four copied scripts use assertions, which must remain enabled.
Their additional historical comparison outputs are not proof premises
of this supplement. Optional floating diagnostics are numerical
observations only. None of these checks certifies a finite Ringmin global
optimum. The full all-pairs and infinite-parameter statements are proved
in the text, including both cyclic arcs, actual-label recovery, all seams,
the root scale before asymptotic substitution, and both parities.

\end{document}
